\documentclass{amsart}
% For dealing with references we use the comment environment
\usepackage{verbatim}
\newenvironment{reference}{\comment}{\endcomment}
%\newenvironment{reference}{}{}
\newenvironment{slogan}{\comment}{\endcomment}
\newenvironment{history}{\comment}{\endcomment}

% For commutative diagrams we use Xy-pic
\usepackage[all]{xy}

% We use 2cell for 2-commutative diagrams.
\xyoption{2cell}
\UseAllTwocells

% We use multicol for the list of chapters between chapters
\usepackage{multicol}

% This is generally recommended for better output
\usepackage{lmodern}
\usepackage[T1]{fontenc}

% For cross-file-references

% Package for hypertext links:
\usepackage{hyperref}

% For any local file, say "hello.tex" you want to link to please

% Theorem environments.
%
\theoremstyle{plain}
\newtheorem{theorem}[subsection]{Theorem}
\newtheorem{proposition}[subsection]{Proposition}
\newtheorem{lemma}[subsection]{Lemma}

\theoremstyle{definition}
\newtheorem{definition}[subsection]{Definition}
\newtheorem{example}[subsection]{Example}
\newtheorem{exercise}[subsection]{Exercise}
\newtheorem{situation}[subsection]{Situation}

\theoremstyle{remark}
\newtheorem{remark}[subsection]{Remark}
\newtheorem{remarks}[subsection]{Remarks}

\numberwithin{equation}{subsection}

% Macros
%
\def\lim{\mathop{\mathrm{lim}}\nolimits}
\def\colim{\mathop{\mathrm{colim}}\nolimits}
\def\Spec{\mathop{\mathrm{Spec}}}
\def\Hom{\mathop{\mathrm{Hom}}\nolimits}
\def\Ext{\mathop{\mathrm{Ext}}\nolimits}
\def\SheafHom{\mathop{\mathcal{H}\!\mathit{om}}\nolimits}
\def\SheafExt{\mathop{\mathcal{E}\!\mathit{xt}}\nolimits}
\def\Sch{\mathit{Sch}}
\def\Mor{\mathop{\mathrm{Mor}}\nolimits}
\def\Ob{\mathop{\mathrm{Ob}}\nolimits}
\def\Sh{\mathop{\mathit{Sh}}\nolimits}
\def\NL{\mathop{N\!L}\nolimits}
\def\CH{\mathop{\mathrm{CH}}\nolimits}
\def\proetale{{pro\text{-}\acute{e}tale}}
\def\etale{{\acute{e}tale}}
\def\QCoh{\mathit{QCoh}}
\def\Ker{\mathop{\mathrm{Ker}}}
\def\Im{\mathop{\mathrm{Im}}}
\def\Coker{\mathop{\mathrm{Coker}}}
\def\Coim{\mathop{\mathrm{Coim}}}

% Boxtimes
%
\DeclareMathSymbol{\boxtimes}{\mathbin}{AMSa}{"02}

%
% Macros for moduli stacks/spaces
%
\def\QCohstack{\mathcal{QC}\!\mathit{oh}}
\def\Cohstack{\mathcal{C}\!\mathit{oh}}
\def\Spacesstack{\mathcal{S}\!\mathit{paces}}
\def\Quotfunctor{\mathrm{Quot}}
\def\Hilbfunctor{\mathrm{Hilb}}
\def\Curvesstack{\mathcal{C}\!\mathit{urves}}
\def\Polarizedstack{\mathcal{P}\!\mathit{olarized}}
\def\Complexesstack{\mathcal{C}\!\mathit{omplexes}}
% \Pic is the operator that assigns to X its picard group, usage \Pic(X)
% \Picardstack_{X/B} denotes the Picard stack of X over B
% \Picardfunctor_{X/B} denotes the Picard functor of X over B
\def\Pic{\mathop{\mathrm{Pic}}\nolimits}
\def\Picardstack{\mathcal{P}\!\mathit{ic}}
\def\Picardfunctor{\mathrm{Pic}}
\def\Deformationcategory{\mathcal{D}\!\mathit{ef}}

\setlength{\emergencystretch}{3em}
\begin{document}
\title[Stacks Project, Tag 0EC2]{313 --- Injectivity of residue-supported modules with characteristic-zero derivation operators}
\maketitle
\noindent Copyright (C) 2005--2025 Johan de Jong. Stacks Project authors. Source: \href{https://stacks.math.columbia.edu/tag/0EC2}{Tag 0EC2}.

\noindent Editorial difficulty note (not author text). Difficulty H2: The proof constructs a socle-preserving embedding into a sum of injective hulls. Leibniz operators force divisibility by a parameter, and an induction on dimension then proves surjectivity on parameter torsion. This is a module classification from differential-operator structure with a complete divisibility argument..

\noindent Editorial provenance note (not author text). History: excerpt from revision \texttt{a0444\allowbreak 6e57e\allowbreak c1fbc\allowbreak 252a8\allowbreak 71afc\allowbreak ec775\allowbreak 2fb28\allowbreak 07b14}, local-cohomology.tex, lines 3790--3855. Reference presentation was adapted; original-author context and core support blocks were retained or added. The mathematical wording is unchanged. Licensed under GNU FDL 1.2 or later; no invariant sections or cover texts. See ../licenses/COPYING. Context blocks reproduce author definitions and supporting results; they are not additional counted problems.

\begin{lemma}
\label{lemma-structure-torsion-D-module-regular}
\begin{reference}
Special case of \href{https://stacks.math.columbia.edu/bibliography}{Bibliography: Lyubeznik (Theorem 2.4)}
\end{reference}
Let $k$ be a field of characteristic $0$. Let $d \geq 1$.
Let $A = k[[x_1, \ldots, x_d]]$ with maximal ideal $\mathfrak m$.
Let $M$ be an $\mathfrak m$-power torsion $A$-module endowed with
additive operators $D_1, \ldots, D_d$ satisfying the leibniz rule
$$
D_i(fz) = \partial_i(f) z + f D_i(z)
$$
for $f \in A$ and $z \in M$. Here $\partial_i$ is
differentiation with respect to $x_i$.
Then $M$ is isomorphic to a direct sum
of copies of the injective hull $E$ of $k$.
\end{lemma}

\begin{proof}
Choose a set $J$ and an isomorphism $M[\mathfrak m] \to \bigoplus_{j \in J} k$.
Since $\bigoplus_{j \in J} E$ is injective
(Dualizing Complexes, Lemma \href{https://stacks.math.columbia.edu/tag/08XV}{lemma sum injective modules (Tag 08XV)})
we can extend this isomorphism to an $A$-module homomorphism
$\varphi : M \to \bigoplus_{j \in J} E$.
We claim that $\varphi$ is an isomorphism, i.e., bijective.

\medskip\noindent
Injective. Let $z \in M$ be nonzero. Since $M$ is $\mathfrak m$-power torsion
we can choose an element $f \in A$ such that $fz \in M[\mathfrak m]$ and
$fz \not = 0$. Then $\varphi(fz) = f\varphi(z)$ is nonzero, hence
$\varphi(z)$ is nonzero.

\medskip\noindent
Surjective. Let $z \in M$. Then $x_1^n z = 0$ for some $n \geq 0$.
We will prove that $z \in x_1M$ by induction on $n$.
If $n = 0$, then $z = 0$ and the result is true.
If $n > 0$, then applying $D_1$ we find $0 = n x_1^{n - 1} z + x_1^nD_1(z)$.
Hence $x_1^{n - 1}(nz + x_1D_1(z)) = 0$. By induction we get
$nz + x_1D_1(z) \in x_1M$. Since $n$ is invertible, we conclude
$z \in x_1M$. Thus we see that $M$ is $x_1$-divisible.
If $\varphi$ is not surjective, then we can choose
$e \in \bigoplus_{j \in J} E$ not in $M$.
Arguing as above we may assume $\mathfrak m e \subset M$,
in particular $x_1 e \in M$. There exists an element
$z_1 \in M$ with $x_1 z_1 = x_1 e$. Hence
$x_1(z_1 - e) = 0$. Replacing $e$ by $e - z_1$
we may assume $e$ is annihilated by $x_1$.
Thus it suffices to prove that
$$
\varphi[x_1] :
M[x_1]
\longrightarrow
\left(\bigoplus\nolimits_{j \in J} E\right)[x_1] =
\bigoplus\nolimits_{j \in J} E[x_1]
$$
is surjective. If $d = 1$, this is true by construction of $\varphi$.
If $d > 1$, then we observe that $E[x_1]$ is the injective hull
of the residue field of $k[[x_2, \ldots, x_d]]$, see
Dualizing Complexes, Lemma \href{https://stacks.math.columbia.edu/tag/08Z2}{lemma quotient (Tag 08Z2)}.
Observe that $M[x_1]$ as a module over $k[[x_2, \ldots, x_d]]$
is $\mathfrak m/(x_1)$-power torsion and comes
equipped with operators $D_2, \ldots, D_d$ satisfying
the displayed Leibniz rule.
Thus by induction on $d$ we conclude that $\varphi[x_1]$
is surjective as desired.
\end{proof}
\end{document}
